\section{Exponential profiles and exclusion of infinite growth}
\label{sec:profiles}

Throughout this Section we are in the nonisotropic, noncircular case of
\Cref{thm:shape-lift}. Thus \(f=\phi'\) and \(q=f''/f\) are holomorphic
on \(D_H=\{s:|\Im s|<H\}\), where \(H\) may be infinite, and
\[
 M(s)=\max\{|f(s)|,|q(s)|^{1/2}\}.
\]
The impact lift is denoted by \(\mathfrak s(y,t)\). Its normalization is
\(\mathfrak s(y,0)=y\), and it is even in \(t\). We shall exclude a
particular pair of rapidly varying field profiles. The point of retaining
both profiles in the statement is that the second can occur even when
the normalized first field tends to zero.

\begin{theorem}[Exponential-profile obstruction]
\label{thm:exponential-obstruction}
There is no sequence of centers \(s_n=x_n+\ii h_n\in D_H\), length scales
\(\ell_n>0\), and field scales \(m_n>0\) with the following properties.
Put \(b_n=\ell_nm_n\) and
\begin{equation}
 \label{profile:scales}
 F_n(S)=\frac{f(s_n+\ell_n S)}{m_n},\qquad
 Q_n(S)=\frac{q(s_n+\ell_n S)}{m_n^2},\qquad
 Y=\frac{y-s_n}{\ell_n},\quad T=m_nt.
\end{equation}
Assume that \(b_n\to\infty\), that \((D_H-s_n)/\ell_n\) exhausts
\(\C\), and that, locally uniformly on \(\C\), either
\begin{equation}
 \label{profile:exponential-pairs}
 (F_n,Q_n)\longrightarrow(c_0e^{-\ii S},0)
 \quad\hbox{or}\quad
 (F_n,Q_n)\longrightarrow(0,c_0^2e^{-2\ii S}),
 \qquad c_0\ne0.
\end{equation}
The following additional assumptions are imposed according to the value
of \(H\).
\begin{enumerate}[label=\textnormal{(\roman*)}]
 \item If \(H<\infty\), then \(h_n\to H\), \(\ell_n\to0\), and
 there are \(\alpha_0\in(1/2,1)\) and
 \(\delta_b\in(0,H/2)\) such that, for every fixed \(L<\infty\),
 there is a constant \(C_L\) for which eventually
 \begin{equation}
  \label{profile:coarser-cone}
  M(s)\le C_Lm_n\left(\frac{\ell_n}{\delta}\right)^{\alpha_0},
  \qquad
  \ell_n\le\delta=h_n-\Im s\le\delta_b,
  \quad |\Re s-x_n|\le L\delta.
 \end{equation}
 The exponent is independent of \(L\); the constant and the eventual
 sequence threshold may depend on \(L\).
 \item If \(H=\infty\), then \(h_n\to\infty\),
 \(0<\ell_n\le C\), and \eqref{profile:coarser-cone} holds with a
 single constant for all horizontal positions and for
 \(\ell_n\le\delta\le h_n\). Moreover,
 \begin{equation}
  \label{profile:entire-log-condition}
  \log(1+h_n/\ell_n)=o(\log m_n).
 \end{equation}
\end{enumerate}
\end{theorem}

The proof has three parts. First we initialize the lift far below the
center height, including at complex parameters. Next we identify its
unique limiting equation and a branch point of its solution. Finally,
we use real-parameter width continuation to make that branch point
incompatible with a single-valued impact lift.

\subsection{Joint deep data}

For a local continuation of the lift define
\[
 S_n(Y,T)=\frac{\mathfrak s(s_n+\ell_nY,T/m_n)-s_n}{\ell_n}.
\]
A continuation is said to be reached vertically if its germs are
obtained along the vertical segment with the horizontal position and
parameter fixed, starting in known lower data. Such germs agree on
overlapping neighborhoods: a finite chain of neighborhoods along the
segment compares any two sufficiently close vertical paths.

\begin{lemma}[Deep passages]
\label{profile:deep-data}
Assume the hypotheses in \Cref{thm:exponential-obstruction}, without
assuming its conclusion. Fix \(\alpha\in(1/2,\alpha_0)\) and
\(T_0<\infty\). There are \(D_0\) and \(C\), independent of the
fixed aperture \(A\), with the following property. For every fixed
\(A<\infty\), and eventually in \(n\), there are joint holomorphic
vertical passages on compact subsets of
\begin{equation}
 \label{profile:deep-domain}
 D=-\Im Y>D_0,\qquad |\Re Y|<AD,\qquad |T|<T_0,
\end{equation}
and
\begin{equation}
 \label{profile:deep-estimate}
 |S_n(Y,T)-Y|\le C|T|^2D^{1-2\alpha}.
\end{equation}
These passages start in the actual lift domain and remain in the shape
strip. Their eventual sequence threshold may depend on \(A\) and on
the compact subset.
\end{lemma}

\begin{proof}
We first make a uniformity observation. With \(D=\delta/\ell_n\),
weakening the exponent in \eqref{profile:coarser-cone} gives
\[
 D^\alpha M/m_n\le C_LD^{-(\alpha_0-\alpha)}.
\]
For each fixed \(L\), this is bounded by an absolute constant beyond a
finite depth threshold depending on \(L\). On the remaining finite
cone, \eqref{profile:exponential-pairs} supplies the same bound
eventually: the limiting value of \(M/m_n\) is
\(|c_0|e^{-D}\), independent of the horizontal position. Consequently
the weakened cone bound can have a constant independent of \(L\),
provided its eventual sequence threshold is allowed to depend on
\(L\). This permits the choice of \(D_0\) before a large aperture.

Suppose first that \(H<\infty\). Choose \(D_0\ge10\), so the
required input depth \(d/10\) is at least \(\ell_n\).
Apply \Cref{thm:cone} and its
reference-height formulation in \Cref{cone:reference-height} with reference
height \(h_n\), rather than \(H\), and with
\[
 d=\ell_nD_0,\qquad r=\frac{2T_0}{m_n},\qquad
 \eps=C'T_0D_0^{-\alpha}.
\]
The cone hypothesis is exactly
\(rM\le\eps(d/\delta)^\alpha\). Since
\(r/d=2T_0/(b_nD_0)\to0\), the fixed ratio allowance in that theorem
can be chosen to be one. Its auxiliary bound satisfies
\[
 B=1+(\eps/r)d^\alpha\le C m_n\ell_n^\alpha.
\]
Here \(m_n\ell_n^\alpha\to\infty\), since \(b_n\to\infty\) and
\(\ell_n\to0\). The physical displacement estimate, divided by
\(\ell_n\), is therefore \eqref{profile:deep-estimate}. Increasing
\(D_0\) makes the required smallness conditions hold. The starting
slab lies in a fixed strict substrip of the original shape strip.
The relative displacement is small compared with the depth below
\(h_n\); in particular the image remains in \(D_H\). This use of an
auxiliary reference height needs no no-patch assertion at that height.
The aperture assertion of \Cref{thm:cone}, together with the preceding
weakening argument, gives the stated independence of \(D_0,C\) from
\(A\).

We provide the initializer needed when \(H=\infty\), because the
starting height is then unbounded. Retain the same \(d,r\), and put
\(B=Cm_n\ell_n^\alpha\). All horizontal ranges may now be used.
Condition \eqref{profile:entire-log-condition} implies
\begin{equation}
 \label{profile:entire-start-smallness}
 r(h_n/d)^\alpha\longrightarrow0,
 \qquad Bh_n^{-\alpha}\longrightarrow\infty.
\end{equation}
Thus the known data near \(\Im y=0\) have uniformly small parameters
over the entire range that will be used.

Consider first \(|\Im t|\ge\gamma|t|\), with a fixed steep cutoff
\(\gamma>0\). Use the upward increment \(t\) or \(-t\), as
appropriate. The normalized chord recurrence and kernel estimates in
\Cref{thm:cone} apply along the segment from the real-row slab to depth
\(\delta\ge d/2\), for
\( |t|\le10r(\delta/d)^\alpha\). Indeed the fast fields are small
there after increasing \(D_0\), while \(|t|/\delta\) is bounded.
Local Cauchy disks have radius a fixed fraction of the depth. Any part
of a thickening slightly below the real row is covered by real
symmetry and the comparable bound above that row. The recurrence
estimates, with \(x\) denoting the current depth, are
\[
 |J-1|\le C|t|^2B^2x^{-2\alpha},\qquad
 |\mathfrak s-y|\le C|t|^2B^2x^{1-2\alpha}.
\]
To see the accumulation directly, the change of normalized step is
\(O(|t|^3B^2x^{-2\alpha-1})\). Summation over rows separated by a
constant multiple of \(|t|\) gives the first estimate, and a second
summation gives the second because \(2\alpha>1\). Initial step and
position errors of order \(|t|^2\) contribute at most
\(C|t|^2(1+h_n)\) to the latter. This is absorbed, uniformly at all
\(\delta\le h_n\), because
\[
 B^2h_n^{1-2\alpha}/(1+h_n)\longrightarrow\infty
\]
by \eqref{profile:entire-start-smallness}. The small relative-error
induction consequently closes. Choosing the number of backward steps
locally constantly, with two starting entries inside the known slab,
gives holomorphic germs. Different choices agree by the exact
recurrence and uniqueness of its nearby root.

It remains to initialize all parameter directions at large depths.
Fix a large constant \(K_0\). At a target depth
\(h_n/K_0\le\delta\le h_n\), set
\(R_t=r(\delta/d)^\alpha\). For \(|t_c|<R_t\) consider
\begin{equation}
 \label{profile:entire-initializer}
 t'=t_c+3R_t\zeta,\qquad
 y=x+\ii\sigma+\sigma g(\zeta),\qquad
 0\le\sigma\le h_n-\delta,
 \quad g(0)=0.
\end{equation}
Choose a smaller fixed steep cutoff for this initializer. On the unit
circle the nonsteep arcs for \(t'\) are uniformly short: the circle
has radius \(3R_t\), its center has modulus at most \(R_t\), and
its intersections with a sufficiently narrow horizontal sector are
transverse. We choose a smooth real boundary function of mean zero
which equals \(-1\) on these arcs and slightly larger collars, is
nowhere below \(-1\), and is at most \(1/(6K_0)\). This is possible
by making the arcs sufficiently short and balancing the negative
bumps by a small positive part. Approximate it, with its mean removed,
by a real trigonometric polynomial with error \(o(1/h_n)\). It is the
boundary imaginary part of a holomorphic polynomial \(g\) with
\(g(0)=0\). In particular,
\[
 1+\Im g=o(1/h_n)\ \hbox{on the nonsteep collars},\qquad
 -o(1/h_n)\le1+\Im g\le1+1/(3K_0)
\]
on the whole circle. The smooth bumps can be chosen with uniform
smooth bounds before approximation, so \(g\) is uniformly bounded;
its degree need not be fixed.

On the nonsteep collars, \eqref{profile:entire-initializer} lands in
the original fixed real-row neighborhood. On the other arcs its depth
is at least
\(\delta-h_n/(3K_0)\ge2\delta/3\), and the steep estimates just
proved apply, or the point is still in the base neighborhood.
All parameter sizes tend to zero by
\eqref{profile:entire-start-smallness}. These boundary germs glue as
vertical continuations from below.

Here and in subsequent disk fillings we use the following explicit
holomorphic continuation argument. Pull the boundary data back to a
thin annulus about \(|\zeta|=1\), times a complex neighborhood of
the center path. Their negative Laurent coefficients in \(\zeta\)
are holomorphic functions of the center variables. Initially the full
disks lie in known data, so all those coefficients vanish. The
identity theorem makes them vanish along the whole connected center
path. The nonnegative Laurent series fills each disk and gives joint
holomorphic center germs. Applied to
\((\mathfrak s-y)/(t')^2\), which is regular at \(t'=0\) in the
initial data, the maximum principle transfers the boundary estimate
to the disk center. Small complex changes of \(\sigma,t_c\) are
allowed by the strict margins above, so these really are joint germs.
This proves the all-direction estimate for
\(h_n/K_0\le\delta\le h_n\), with constants independent of
\(h_n\).

The dyadic disk construction in \Cref{thm:cone} now continues it from
\(h_n/K_0\) to \(d\). Choose \(K_0\) first so large that its
previous-band dips and starting slabs remain at positive heights.
No factor accumulates with the number of bands: the exponent
\(1-2\alpha<0\) makes the estimates on the deeper, previously known
arcs strictly smaller, whereas the steep-arc estimates have an
independent fixed constant. Fix both steep cutoffs and \(K_0\), and
only then increase \(D_0\) to meet their smallness requirements.
After scaling by \(\ell_n\), the result is
\eqref{profile:deep-estimate}. In the infinite-width case the shape
constraint is automatic. This completes both cases.
\end{proof}

\subsection{The unique limiting energy equation}

We use the normalized chord action from
\Cref{lem:kernel-division}. Its frozen versions, with the quotients
continued at zero, are
\begin{align}
 \label{profile:frozen-actions}
 p_F(\lambda,J)&=\frac{J}{\lambda^2}
             -\frac{2\sin(\lambda J/2)}{\lambda^3},&
 E_F(\lambda,J)&=
 \frac{2(\sin z-z\cos z)}{\lambda^3},\\
 p_Q(\lambda,J)&=\frac{J}{2\lambda^2}
             -\frac{\tanh(\lambda J/2)}{\lambda^3},&
 E_Q(\lambda,J)&=
 \frac{\tanh z-z\operatorname{sech}^2z}{\lambda^3},
 \qquad z=\lambda J/2.\notag
\end{align}
Here \(E=Jp_J-p\). In the first case the frozen fast fields are
\((\lambda,0)\), and in the second they are \((0,\lambda^2)\).

\begin{lemma}[Deep model and its uniqueness]
\label{profile:energy-model}
On the deep domains of \Cref{profile:deep-data}, the passages converge
locally uniformly, with derivatives, to a unique holomorphic model
\(S_*(Y,T)\). With \(\lambda=Tc_0e^{-\ii S_*}\) and
\(J=\partial_Y S_*\), it satisfies
\begin{equation}
 \label{profile:energy-equation}
 E_F(\lambda,J)=1/12
 \quad\hbox{or}\quad E_Q(\lambda,J)=1/12,
\end{equation}
respectively, and \(S_*-Y\to0\), \(J\to1\) as
\(-\Im Y\to\infty\) in each fixed cone.
\end{lemma}

\begin{proof}
The deep estimates and Cauchy's inequalities give locally normal
families and the asserted anchors for any subsequential limit. Fix
\(T\ne0\) and put \(a_n=T/b_n\). On a compact set with a deep
margin,
\[
 J_n(Y)=\frac{S_n(Y+a_n,T)-S_n(Y,T)}{a_n}
       \longrightarrow\partial_Y S_*(Y,T).
\]
The fast fields used in the normalized action are
\(TF_n(S+a_nu)\) and \(T^2Q_n(S+a_nu)\). They converge analytically
to the appropriate frozen pair in
\eqref{profile:frozen-actions}. The action \(p_n(S,J;T)\) and its
derivatives converge on these neighborhoods by
\Cref{cone:profile-action-convergence}. In particular the case
\(F_n\to0\) uses the analytic division in
\Cref{lem:kernel-division}, rather than division of two limiting
zeros.

The exact recurrence takes the form
\begin{equation}
 \label{profile:scaled-recurrence}
 p_{n,J}(S,J)-a_np_{n,S}(S,J)
   =p_{n,J}(S_-,J_-),\qquad S=S_-+a_nJ_-.
\end{equation}
It holds on the passages by analytic continuation. To fix its root
choice, first take \(T\) sufficiently small, for each fixed \(n\),
so that all physical points in question lie in known strict
substrips with the glancing branch. Continue in the deep \(T\)-disk.
Its shifted arguments remain in the domain with margin, and Cauchy
estimates keep \(J_n\) near one. The fast fields are uniformly small
after increasing \(D_0\), so the relevant kernels and roots remain
nondegenerate throughout that continuation.

Divide the difference in \eqref{profile:scaled-recurrence} by
\(a_n\). Both terms are evaluations of the same analytic function
\(p_{n,J}\) at neighboring states. Analytic convergence therefore
gives
\[
 \frac{\dd}{\dd Y}p_J(S_*,J)=p_S(S_*,J).
\]
This step requires no rate of convergence compared with \(a_n\).
Since \(S_{*,Y}=J\), differentiating \(Jp_J-p\) shows that it is
constant in \(Y\). The deep anchor gives its value \(1/12\).
All statements hold on arbitrary compact sets by exhaustion, and at
\(T=0\) by removability.

For uniqueness write \(w=Tc_0e^{-\ii Y}\) and
\(z=\lambda J/2\). The energy equation is
\begin{equation}
 \label{profile:lambda-equations}
 \lambda^3=24(\sin z-z\cos z)
 \quad\hbox{or}\quad
 \lambda^3=12(\tanh z-z\operatorname{sech}^2z).
\end{equation}
At \((\lambda,J)=(0,1)\), \(\partial_J E=1/4\); thus it determines
\(J\) analytically as a function of \(\lambda\) near that point.
Moreover
\[
 w\frac{\dd\lambda}{\dd w}=\lambda J=2z,
 \qquad \frac{\lambda}{w}\longrightarrow1.
\]
Integrating this differential relation gives
\begin{equation}
 \label{profile:inverse-lambda}
 w=\lambda\exp\left(\int_0^\lambda
                  \frac{1/J(v)-1}{v}\,\dd v\right),
 \qquad S_*-Y=\ii\log(\lambda/w).
\end{equation}
The integrand is analytic at zero, the inverse has derivative one
there, and the anchor fixes the constant and logarithm. Hence there
is a unique deep model. Normality now gives convergence of the whole
sequence on its compact domains, with derivatives by Cauchy's
inequalities.
\end{proof}

\subsection{A branch point of the model}

\begin{lemma}[Regular model and nontrivial continuation]
\label{profile:branch-model}
The model of \Cref{profile:energy-model}, expressed in the variable
\(w=Tc_0e^{-\ii Y}\), is analytic for \(|w|<R\), for some
\(R>0\). It continues through the boundary circle except possibly
at \(w=R,-R\). A sufficiently small continuation loop around
\(w=R\), attached to the inner model, changes the value of
\(S_*-Y\). The loop and its two adjoining regular corridors can be
chosen so that all chord kernels and the twist \(p_{JJ}\) remain
nonzero, including on the changed branch.
\end{lemma}

\begin{proof}
In the first case of \eqref{profile:lambda-equations} put \(v=z\)
and define
\begin{equation}
 \label{profile:trigonometric-g}
 g(v)=\frac{3(\sin v-v\cos v)}{v^3},\qquad
 P(v)=\frac{v g(v)^{2/3}}{\sin v},\qquad v_c=\pi.
\end{equation}
In the second case put \(v=\tanh z\) and define
\begin{equation}
 \label{profile:hyperbolic-g}
 g(v)=\frac{3}{2v^3}
       \bigl[v-(1-v^2)\operatorname{arctanh}v\bigr],\qquad
 P(v)=g(v)^{2/3},\qquad v_c=1.
\end{equation}
All branches at zero are normalized by \(g(0)=P(0)=1\).
Direct differentiation of \eqref{profile:lambda-equations} gives,
in both cases,
\begin{equation}
 \label{profile:v-equation}
 \lambda=2v g(v)^{1/3},\qquad
 wv'(w)=vP(v),\qquad v(w)\sim w/2.
\end{equation}
For example, in the first case
\(\lambda_v=2\sin v/(vg(v)^{2/3})\); in the second it is
\(2\operatorname{arctanh}v/(vg(v)^{2/3})\).

We verify the positivity and zero-freeness needed to control the
analytic region. In \eqref{profile:trigonometric-g}, every nonzero
zero of \(\sin v-v\cos v\) is real. Indeed
\(u(x)=\sin(vx)\) solves
\(-u''=v^2u\), \(u(0)=0\), \(u'(1)=u(1)\), and integration by
parts gives the nonnegative quadratic form
\[
 \int_0^1|u'|^2\,\dd x-|u(1)|^2\ge0;
\]
the inequality follows from \(u(1)=\int_0^1u'\). Hence \(v^2\) is
real and nonnegative. The positive roots \(r_k\) are simple and
satisfy \(k\pi<r_k<k\pi+\pi/2\), by \(\tan v=v\).
Hadamard factorization for the even entire function of order one
\cite[Theorem~3.7, p.~43]{McMullen2017} therefore gives
\[
 g(v)=\prod_{k\ge1}(1-v^2/r_k^2).
\]
Comparison with the sine product yields
\[
 \log P(v)=\sum_{j\ge1}\frac{v^{2j}}{j}
 \sum_{k\ge1}\left((k\pi)^{-2j}-\frac23r_k^{-2j}\right).
\]
All its coefficients are strictly positive, and it is analytic for
\(|v|<\pi\). In particular \(P\) has nonnegative even
coefficients, with a positive quadratic coefficient.

In \eqref{profile:hyperbolic-g},
\[
 g(v)=\sum_{j\ge0}a_jv^{2j},\qquad
 a_j=\frac{3}{(2j+1)(2j+3)},\qquad a_0=1.
\]
Write its reciprocal formally as
\(1-\sum_{j\ge1}b_jv^{2j}\). The ratios
\(a_j/a_{j-1}=(2j-1)/(2j+3)\) increase. The recurrence
\[
 b_n=a_n-\sum_{j=1}^{n-1}b_ja_{n-j}
\]
shows by induction that \(0\le b_n\le a_n\). For completeness,
with \(r_n=a_n/a_{n-1}\), the induction hypothesis gives
\[
 \sum_{j=1}^{n-2}b_ja_{n-j}
 \le r_n(a_{n-1}-b_{n-1}),\qquad
 b_n\ge b_{n-1}(r_n-a_1)\ge0.
\]
These series converge in the unit disk. For real \(0\le r<1\),
\(\sum b_jr^{2j}=1-1/g(r)<1\); thus its modulus is less than one
throughout that disk. Consequently \(g\) is zero-free there and
\(P=(1-\sum b_jv^{2j})^{-2/3}\) has nonnegative even
coefficients, again with a positive quadratic coefficient.

The inverse of \eqref{profile:v-equation} is analytic at zero:
\[
 w(v)=2v\exp\left(\int_0^v
                    \frac{1/P(\xi)-1}{\xi}\,\dd\xi\right).
\]
The differential recurrence for its inverse shows that \(v(w)\) is
odd and has nonnegative Taylor coefficients. Its linear coefficient
is \(1/2\), and its cubic coefficient is strictly positive. On
\(0<v<v_c\), \(w(v)\) increases and
\[
 \frac{\dd\log w}{\dd v}=\frac{1}{vP(v)}.
\]
It tends to a finite positive limit \(R\) as \(v\uparrow v_c\):
the last integrand is integrable at \(\pi\) in the first case and
has a finite limit at one in the second.

The power series of \(v(w)\) reaches radius \(R\). Otherwise, at
a smaller radius \(\rho\), positivity bounds its sum of absolute
coefficients by the real solution \(v(\rho)<v_c\). It converges
uniformly on the closed circle and the regular differential equation
\eqref{profile:v-equation} continues it through every point of that
circle, a contradiction. At radius \(R\) its positive sum is
\(v_c\). At phases other than zero and \(\pi\), the modulus is
strictly smaller than \(v_c\): equality in the triangle inequality
for the strictly positive linear and cubic terms would require
\(e^{2\ii\theta}=1\). The differential equation therefore gives
continuation through those phases.

For \(0<|w|<R\), the solution \(v(w)\) never vanishes, since
uniqueness for the regular differential equation at a nonzero \(w\)
would otherwise make it identically zero. It follows that
\(\lambda/w\) is nonzero and holomorphic on the whole disk, with
value one at zero. Thus \eqref{profile:inverse-lambda} defines the
impact model there. Its regularity tests can be made explicit. In
the first case,
\[
 J=g(v)^{-1/3},\qquad p_{JJ}=\frac{\sin v}{2\lambda},\qquad
 k_0=J,\quad k_+=k_-=\frac{2\sin v}{\lambda}.
\]
In the second case, with \(z=\operatorname{arctanh}v\),
\[
 J=\frac{2z}{\lambda},\qquad
 p_{JJ}=\frac{v(1-v^2)}{2\lambda},\qquad
 k_0=k_+=k_-=\frac{2v}{(1-v^2)\lambda}.
\]
They are nonzero in the inner disk, with removable values
\(J=k_j=1\), \(p_{JJ}=1/4\) at zero.

We finally verify that continuation at \(R\) actually changes the
impact model. In the first case \(w(v)\) is analytic at \(v=\pi\)
and has a simple quadratic critical point there. If
\(\lambda_c=(24\pi)^{1/3}\), then
\(w''(\pi)=-4R/\lambda_c^2\ne0\). Hence \(v(w)\) has a
nonzero square-root term. The identity \(w\lambda_w=2v\) shows
that \(\lambda(w)\), though its square-root term cancels, has a
nonzero term of order \((w-R)^{3/2}\). More explicitly, on the
branch from \(0<w<R\), with \(\eta=1-w/R\),
\[
 \lambda=\lambda_c-2\pi\eta
        +\frac{2\sqrt2\lambda_c}{3}\eta^{3/2}+O(\eta^2).
\]
A turn around \(R\) changes this value.

In the second case let \(x=v-1\), \(G=3/2\), and
\(\Lambda=2G^{1/3}\). On any fixed bounded range of turning on the
logarithmic cover,
\[
 g(v)=G\bigl[1-x\log(-x)+(\log2-2)x
              +O(x^2(1+|\log|x||))\bigr].
\]
Integration of \(w_v=w/(vg(v)^{2/3})\) gives
\begin{equation}
 \label{profile:logarithmic-inverse}
 w(v)=R+a_0x+O(x^2(1+|\log|x||)),\qquad
 a_0=R/G^{2/3}\ne0.
\end{equation}
The inverse on a lifted sector exists with
\(x=(w-R)/a_0+O(|w-R|^2(1+|\log|w-R||))\). For example,
Rouche's theorem on disks of a sufficiently small fixed relative
radius gives this inverse and its continuation; the derivative
tends to \(a_0\). Substitution in \(\lambda=2vg(v)^{1/3}\) gives,
with \(\delta=w-R\),
\[
 \lambda=\Lambda-\frac{\delta}{R}\log(-\delta/a_0)
       +\frac{1+\log2}{R}\delta
       +O\bigl(|\delta|^2(1+|\log|\delta||^2)\bigr).
\]
The same coefficients apply over any fixed bounded amount of turning.
The logarithmic coefficient is nonzero, so a turn changes the value;
the remainder is smaller than that change on a sufficiently small
fixed bypass.

These branches attach to the inner model along the real interval
\(w<R\). On a sufficiently small compact bypass, in the first case
\(v\) stays near but not at \(\pi\). In the second it stays near
but not at one, with \(v\ne0,\pm1\), and the chosen
\(\operatorname{arctanh}v\) is finite and nonzero on each compact
lifted corridor. The displayed kernel and twist formulas are
therefore finite and nonzero, also after the turn. Since \(\lambda\)
changes and remains nonzero, \(\ii\log(\lambda/w)\) changes as
well. This proves the Lemma.
\end{proof}

\subsection{Tracking and passage through regular corridors}

\begin{lemma}[Tracking a regular model corridor]
\label{profile:tracking}
Let a compact corridor of a continued branch of the model in
\Cref{profile:branch-model} have nonzero kernels and twist, and let
\(T\) range over a compact set avoiding zero. Suppose two neighboring
lift entries are initialized in a joint patch where they converge to
that model, with convergence of their first divided difference.
Along segments of bounded length made of increments
\(a_n=T/b_n\), the normalized chord recurrence continues these
entries holomorphically and tracks the model with error \(o(1)\).
The same holds for increments \(-T/b_n\). The statement is uniform
on compact thickenings of such segments and their initial patches.
\end{lemma}

\begin{proof}
Cover the compact model states by finitely many analytic action
charts with kernels and \(p_{JJ}\) bounded away from zero. The
analytic action convergence in
\Cref{cone:profile-action-convergence} applies on these charts,
following the branches initialized in the known data. Given the old
state \((S_-,J_-)\), set \(S=S_-+a_nJ_-\). Equation
\eqref{profile:scaled-recurrence} has a unique nearby root for \(J\)
by the implicit function theorem. Evaluating its left side at
\(J=J_-\), and subtracting before taking the limit, gives
\[
 J-J_-=a_n\frac{p_S-J_-p_{SJ}}{p_{JJ}}(S_-,J_-)+o(|a_n|).
\]
The error is uniform with state derivatives on the compact charts.
Indeed the corresponding expression with \(p_n\) has a Taylor
remainder \(O(|a_n|^2)\), and its coefficient converges with
derivatives to the displayed one. The state map consequently is
\begin{equation}
 \label{profile:tracking-map}
 (S_-,J_-)\longmapsto(S_-,J_-)
 +a_n\left(J_-,\frac{p_S-J_-p_{SJ}}{p_{JJ}}(S_-,J_-)\right)
 +o(|a_n|),
\end{equation}
and has local Lipschitz bound \(1+C|a_n|\).

The model solves the ordinary differential equation in
\eqref{profile:tracking-map}: its second component follows by
expanding \(\dd_Yp_J=p_S\). In \(O(1/|a_n|)\) steps, the initial
error and the local \(o(|a_n|)\) errors sum, with a bounded stability
factor, to \(o(1)\). This proves tracking and keeps the continued
states in the chosen regular charts. The shape arguments remain in
the exhausting shape domains. The implicit roots are holomorphic
also in the initial point and parameter; the same estimates on
slightly larger compact charts give uniform convergence on
thickenings. For the negative increment use the fast fields with
\(-T\). The frozen actions are even in \(\lambda\), so their limit
is unchanged, and the same calculation applies.

At an initialization obtained by previous regular passage, the
required exact recurrence and its branch choice persist by analytic
identity from the deep data. Cauchy's integral formula gives
convergence of the first divided differences whenever the two-entry
functions converge on neighborhoods. Thus the construction can be
reinitialized in every regular patch it reaches.
\end{proof}

One consequence will be used repeatedly. If
\(|\Im T|\ge\eta>0\), choose the upward increment
\(a_n=\pm T/b_n\). A target can be reached from an initialized
lower slab whenever its backward slanted segment remains in one
regular corridor. Choose the number of backward steps locally
constantly so that two entries lie well inside the slab, and then
apply \Cref{profile:tracking}. The germs from different such choices
agree: within the initial slab their overlapping entries obey the
exact recurrence, and every subsequent nearby root is unique.
Doing this at intermediate targets gives the germs along the
vertical segment from below. There is room for small changes in both
the target and the parameter, so the passage is joint and converges
on compact neighborhoods. Its horizontal drift is bounded by a
constant times the vertical distance divided by \(\eta\). Thus any
fixed positive \(\eta\), however small, requires only a fixed
finite horizontal range before \(n\) tends to infinity.

\begin{proof}[Proof of \Cref{thm:exponential-obstruction}]
Let
\[
 T_* = |c_0|^{-1},\qquad r_0=T_*/32,
\]
so that \(T_*\) is real and \(|T_*c_0|=1\). Take the parameter
bound in \Cref{profile:deep-data} larger than all the parameter
disks below. We first reach, vertically and with joint convergence,
the inner regular region
\begin{equation}
 \label{profile:inner-region}
 |Tc_0e^{-\ii Y}|\le Re^{-\nu},\qquad
 |T-T_*|\le2r_0,
\end{equation}
for any fixed \(\nu>0\) and every needed compact horizontal range.

Fix a center \(T_c\) in this parameter disk and use
\begin{equation}
 \label{profile:inner-disks}
 T'=T_c+8r_0\zeta,\qquad
 Y'=Y_c-\ii\log(T'/T_c)+g(\zeta),\qquad g(0)=0.
\end{equation}
The parameter disk is bounded away from zero, so this logarithm is
single-valued. Its sign is important:
\[
 T'e^{-\ii[Y_c-\ii\log(T'/T_c)]}=T_ce^{-\ii Y_c}.
\]
Choose a small fixed \(\eta>0\). On the unit circle require the
boundary imaginary part of \(g\) to be a sufficiently large
negative constant on the arcs \(|\Im T'|<\eta\) and slightly
larger collars, and to be at most \(\nu/2\) everywhere. It must
also have mean zero. Such a choice is possible: the bad arcs have
arbitrarily small total length as \(\eta\downarrow0\), uniformly
in the allowed centers, by transversality of the parameter circles.
Use smooth negative bumps on those arcs, balance their mean by a
positive part smaller than \(\nu/2\), then approximate by a
trigonometric polynomial with strict margins. The corresponding
holomorphic polynomial can be normalized by \(g(0)=0\).

On the bad arcs the adjusted point lies in deep initialized data.
On the other arcs, use the upward increment \(\pm T'/b_n\) to
track from a deep slab. Every intermediate point of its backward
slanted segment has smaller \(|T'c_0e^{-\ii Y}|\) than the target,
so the whole segment lies in the inner regular model region. The
horizontal range needed for these paths, including the shift by
\(g\), can be large, but it is fixed after \(\eta\) is chosen.
The depth \(D_0\) in \Cref{profile:deep-data} was independent of
that aperture, so the initialization remains available. The boundary
germs agree as vertical continuations and converge to the model on
compact neighborhoods.

Raise \(\Im Y_c\) from a deep initial position to the desired
target in \eqref{profile:inner-region}. By the maximum principle for
\(\Im g\), the full model disks in \eqref{profile:inner-disks}
stay within radius \(Re^{-\nu/2}\). Initially the full disks lie
in known data. The Laurent-coefficient argument used in the proof of
\Cref{profile:deep-data} fills the disks as the center rises, giving
joint center germs. Apply the maximum principle to the difference
between the filled lift and the analytic model. Boundary convergence
gives convergence at the centers, and on neighborhoods by retaining
the strict margins. This proves \eqref{profile:inner-region}. The
exact chord identities, including neighboring arguments, continue
through these passages for all sufficiently large \(n\).

We next reach two regular patches just beyond the model circle,
near the phases
\begin{equation}
 \label{profile:outer-patches}
 w=Re^{u+\ii\theta}\quad\hbox{and}\quad
 w=Re^{u-\ii\theta},\qquad u\le\nu,
 \quad T=T_*.
\end{equation}
First fix a sufficiently small \(\theta>0\); then choose
\(\nu>0\) sufficiently small in terms of \(\theta\). Use
parameter circles of radius \(r_0\), so that their initial points
remain in the parameter range already reached, and use the same
logarithmic adjustment as in \eqref{profile:inner-disks}. This time
choose a polynomial with
\[
 \Im g\le-4\nu\ \hbox{on }|\Im T'|<\eta_1
 \hbox{ and collars},\qquad
 \Im g\le\nu,\qquad |g|\le C_1\nu.
\]
The cutoff \(\eta_1>0\) and constant \(C_1\) can be fixed
independently of \(\nu\): first choose the sufficiently short bad
arcs, construct one mean-zero adjustment, and multiply it by
\(\nu\).

Raise the centers from \(u=-C\nu\), with \(C\) a sufficiently
large fixed constant, to \(u=\nu\). The nonsteep arcs remain in
the known inner region. On the steep arcs, tracking starts just
below \(u=-\nu\), and the additional vertical distance is
\(O(\nu)\). The horizontal drift is therefore
\(O(\nu/\eta_1)\). Choose \(\nu\) so small that this drift and
the real displacement of \(g\) stay in the chosen regular phase
corridor about each sign in \eqref{profile:outer-patches}. The full
model disks stay in that corridor as well. Disk filling and maximum
modulus, exactly as above, give vertical joint passages and compact
convergence to the two radial model continuations.

At the real physical parameter \(t=T_*/m_n\), one of these passages
reaches the physical height
\[
 h_n+\ell_n(\log R+\nu).
\]
It remains in the shape strip: this was proved in the deep part and
follows on the remaining compact path from model convergence and
the exhaustion of scaled shape domains. Consequently it gives
joint meromorphic germs for
\(X=\Gamma(\mathfrak s)\), attached to the primary branch from
below. If the real-parameter joint width ended at an earlier height,
the passage would give a prohibited joint patch at its first wall,
contrary to \Cref{thm:no-patch}. The same argument excludes equality
at the endpoint, because the passage has a neighborhood there.
Thus, for all sufficiently large \(n\),
\begin{equation}
 \label{profile:real-width-passage}
 h_X(T_*/m_n)>h_n+\ell_n(\log R+\nu).
\end{equation}
This is the one place where a local passage is promoted to an entire
horizontal band, and the parameter here is real. By
\Cref{thm:shape-lift}, \(S_n(Y,T_*)\) is a single-valued
holomorphic function on a neighborhood of the full row
\(\Im Y=\log R+\nu/2\).
The two vertical passages agree with this actual lift, by analytic
identity from their lower starting segments.

Follow this row from phase \(+\theta\) to phase \(-\theta\),
increasing \(\Re Y\), with real steps
\(a_n=T_*/b_n>0\). Initialize in the first of the patches
\eqref{profile:outer-patches}. The continued model follows the
outside of \(w=R\). At its endpoint it differs from the independent
radial continuation to the other patch: close the comparison path
through the inner disk, and the resulting loop winds once around
\(R\), while avoiding \(-R\). This is the nontrivial continuation
proved in \Cref{profile:branch-model}. All states of the two paths
and their small thickenings are regular if \(\theta,\nu\) have
been chosen small enough.

For clarity, we verify why recurrence tracking on this horizontal
path must agree with the actual single-valued lift. Take a small
fixed complex ball \(B_*\) in the initial patch, and run the
tracking construction for each initial offset in \(B_*\). Its
successive output germs live on \(B_*+ja_n\). For an offset \(Y\)
and the offset \(Y+a_n\), their first two common entries coincide
by the actual data and the exact recurrence in the initial patch;
take \(B_*\) with room for two extra steps. Uniqueness of every
subsequent nearby root makes the chains coincide on their common
entries. Hence the germs on adjacent translated balls agree on
overlap. They glue, and by the identity theorem agree throughout
with the actual holomorphic function on the band supplied by
\eqref{profile:real-width-passage}. Tracking gives convergence to
the horizontally continued model on these balls.

Choose a number of steps whose endpoint tends to a point in the
second regular patch; the remaining offset tends to zero. Compact
convergence on neighborhoods makes this harmless. At that point
the same actual function would converge both to the horizontally
continued model and to the distinct radial model value. This
contradiction proves the Theorem.
\end{proof}

\subsection{A log-deficit estimate}

We next produce an exponential profile from either an entire
nonconstant first field or growth of infinite order at a finite
wall. The quotient condition \(q=f''/f\) being holomorphic is used
to control the second field from a logarithmic bound on the first.

\begin{lemma}[Log-deficit estimate]
\label{lem:log-deficit}
There is an absolute constant \(C\) such that the following holds.
Let \(f\) be holomorphic and not identically zero on the unit disk,
suppose \(\log|f|\le K\), and suppose \(f''/f\) extends
holomorphically through every zero. Then
\begin{equation}
 \label{profile:log-deficit-bound}
 \sup_{|z|<1/4}|f''(z)/f(z)|^{1/2}
 \le C\left(1+\int_{|z|<1}(K-\log|f(z)|)\,\dd A(z)\right).
\end{equation}
On a disk of radius \(\Delta\), the corresponding left side is
\(\Delta\sup|q|^{1/2}\) and the integral on the right is replaced
by the disk average of the log deficit.
\end{lemma}

\begin{proof}
Put \(U=K-\log|f|\ge0\) and
\(A=1+\int_{\mathbb D}U\). We may assume that this integral is
finite. The distribution \(-\Delta U\) is \(2\pi\) times the
zero-counting measure, including multiplicities. Testing it against
a nonnegative smooth cutoff equal to one on \(|z|\le0.8\) and
supported in the unit disk bounds the number of those zeros by
\(CA\). If they are \(z_j\), listed with multiplicity, then
\[
 V=U+\sum_j\log|z-z_j|
\]
is harmonic on \(|z|<0.8\). On a slightly smaller disk it has
\(L^1\) norm \(O(A)\), since each logarithm has uniformly bounded
\(L^1\) norm there. Interior harmonic estimates bound its first
and second derivatives by \(CA\), with a further fixed margin.

Equivalently, divide \(f\) by its listed zero factors. The
zero-free quotient has a holomorphic logarithm in the disk, whose
first two derivatives are controlled by the derivatives of \(V\).
Off the zeros, differentiating the logarithm of \(f\) gives
\[
 |f''/f|^{1/2}\le CA+C\sum_j|z-z_j|^{-1}
\]
on a fixed smaller disk. Here the second derivative of the
logarithm contributes square roots of sums of
\(|z-z_j|^{-2}\), bounded by the displayed sum. Each summand is
locally integrable in two real dimensions with a uniform integral
bound, and there are at most \(CA\) of them. The resulting
\(L^1\) bound for \(|q|^{1/2}\), followed by submean for that
subharmonic function, proves
\eqref{profile:log-deficit-bound}. Its subharmonicity across the
zeros uses the assumed holomorphic extension of \(q\). Rescale the
variable to obtain the final assertion.
\end{proof}

We record the periodic form that will also be useful below. The
horizontal mean
\[
 \mathcal L(h)=\frac1{2\pi}\int_0^{2\pi}
                     \log|f(x+\ii h)|\,\dd x
\]
is convex in height by subharmonicity, and even by real symmetry.
Thus \(\mathcal L(h)\ge\mathcal L(0)>-\infty\); the real-row
value is finite since \(f>0\) there. If \(K\) bounds
\(\log|f|\) at all heights of disks of radius \(\Delta\),
averaging their average deficits over the horizontal centers gives
\(O(1+K_+)\), with constants depending only on the fixed real-row
data. For each individual center positivity and periodicity give
the bound
\begin{equation}
 \label{profile:periodic-deficit}
 1+\text{disk-average}(K-\log|f|)
 \le C(1+\Delta^{-1})(1+K_+).
\end{equation}
Indeed, enclose the disk in a rectangle of width and height
\(2\Delta\). At each height an interval of that width is covered
by \(O(1+\Delta)\) whole periods; its integral of the nonnegative
deficit is at most \(C(1+\Delta)(1+K_+)\). Integration in height
and division by the disk area gives
\eqref{profile:periodic-deficit}. In particular, on a smaller disk,
\begin{equation}
 \label{profile:periodic-q-bound}
 |q|^{1/2}\le C\Delta^{-1}(1+\Delta^{-1})(1+K_+).
\end{equation}

\subsection{Selecting a profile at entire or infinite-order growth}

\begin{theorem}[Exclusion of entire and infinite-order growth]
\label{thm:infinite-growth-exclusion}
In the nonisotropic case, \(H=\infty\) implies that the table is a
circle. If \(H<\infty\), then
\[
 \limsup_{k\to\infty}
 \frac{\log\max_x M(x+\ii(H-e^{-k}))}{k}<\infty.
\]
\end{theorem}

\begin{proof}
We may exclude constant \(f\), which gives a circle. Suppose first
that \(H<\infty\) and the displayed limsup is infinite. The same
is true with \(M\) replaced by \(|f|\). Otherwise
\(\log\max|f|=O(k)\) on those rows; disks with radius comparable
to \(e^{-k}\), and with a fixed extra depth margin, have a global
log upper bound \(O(k)\). Formula
\eqref{profile:periodic-q-bound} then gives a finite power bound in
inverse depth for \(|q|^{1/2}\), contradicting the assumed infinite
order of \(M\).

Treat this case and \(H=\infty\) together. Put
\[
 \rho(h)=\log\max_x|f(x+\ii h)|.
\]
The function \(\rho\) is convex in height by the three-lines
theorem, and even by symmetry. In the finite-width case just
identified it tends to infinity and has infinite upper order
relative to \(\log(1/(H-h))\). In the entire case a nonconstant
real-symmetric periodic function has a growing Fourier mode:
if a nonconstant coefficient is nonzero, its conjugate coefficient
is also nonzero, and one grows as \(h\to\infty\). Hence
\(\rho(h)\ge c h-C\) for some \(c>0\). In either case
\(\rho\to\infty\), and it is eventually strictly increasing
with slope bounded below by a positive constant.

Let \(h(R)\) denote its inverse at large values, let \(N(R)\) be
its right slope there, and put \(\mu(R)=\log N(R)\). The function
\(\mu\) is nondecreasing and bounded below. The inverse is locally
absolutely continuous and
\[
 h'(R)=e^{-\mu(R)}\quad\hbox{for almost every }R.
\]
We claim that
\begin{equation}
 \label{profile:slow-slope-values}
 \liminf_{R\to\infty}\mu(R)/R=0.
\end{equation}
If \(\mu(R)\ge cR\) eventually for some \(c>0\), inverse
integration would give a finite limiting height and
\(H-h(R)\le Ce^{-cR}\). This is impossible if \(H=\infty\).
For finite \(H\) it implies
\(R/\log(1/(H-h(R)))\le1/c+o(1)\), contradicting infinite
upper order. Since \(\mu\) is bounded below, the liminf cannot be
negative. This proves \eqref{profile:slow-slope-values}.

Choose \(R\to\infty\) along values with \(\mu(R)=o(R)\), and
fix a lower base value \(R_b\). The positive measure \(\dd\mu\)
on \([R_b,R]\) has mass \(o(R)\). The weak maximal-interval
inequality gives a number \(\eps_R\to0\) and a continuity point
\(v\in[R/3,2R/3]\) such that
\begin{equation}
 \label{profile:flat-slope-selection}
 |\mu(z)-\mu(v)|\le\eps_R|z-v|
 \qquad (R_b\le z\le R).
\end{equation}
Here is a direct measure justification of the choice. The set of
points belonging to an interval on which the mass divided by its
length exceeds \(\eps_R\) has length at most a constant times
the total mass divided by \(\eps_R\), by the disjoint-interval
covering argument. Choose, for example, a constant multiple of
\(\sqrt{(1+\mu(R)-\mu(R_b))/R}\) for \(\eps_R\). The bad set
has length \(o(R)\), so cannot cover \([R/3,2R/3]\).
Increasing the constant accounts for intervals with one endpoint
at the selected point. Atoms and discontinuity points may be
discarded. This gives both one-sided secant bounds in
\eqref{profile:flat-slope-selection}.

Set
\begin{equation}
 \label{profile:selected-scales}
 h_n=h(v),\qquad \ell_n=N(v)^{-1},\qquad m_n=e^v.
\end{equation}
For bounded \(u\), integrate the inverse derivative using
\eqref{profile:flat-slope-selection}, then invert the result. It
gives, uniformly on each fixed bounded real range,
\begin{equation}
 \label{profile:linear-log-maximum}
 \rho(h_n+\ell_nu)=v+u+o(1).
\end{equation}
These heights stay in the shape strip: bounded intervals of values
about \(v\) lie in \([R_b,R]\) with distances tending to infinity
from the endpoints, and the inverse integration applies there.
In particular the upper-wall distance in units of \(\ell_n\)
tends to infinity.

The slope lower bound, and \(\mu(v)\le\mu(R)=o(R)\), also give
\begin{equation}
 \label{profile:selected-size-bounds}
 0<\ell_n\le C,\qquad h_n\le Cv,
 \qquad |\log\ell_n|=o(v),\qquad
 b_n=\ell_nm_n=\exp(v-o(v))\longrightarrow\infty.
\end{equation}
The centers satisfy \(h_n\to H\), with the infinite value
allowed. For finite \(H\), convex divergence at that height forces
the slopes to infinity, so \(\ell_n\to0\). For infinite \(H\),
\(h_n\to\infty\) and \(\ell_n\le C\). Thus the lower wall also
escapes in scaled coordinates, and the shape domains exhaust
\(\C\).

Choose \(x_n\) realizing the row maximum at \(h_n\). By
\eqref{profile:linear-log-maximum}, the entire subsequential limits
of
\[
 F_n(S)=m_n^{-1}f(x_n+\ii h_n+\ell_nS)
\]
satisfy \(|F(S)|\le e^{\Im S}\) and \(|F(0)|=1\). The maximum
principle applied to \(e^{\ii S}F(S)\) therefore gives
\begin{equation}
 \label{profile:selected-exponential}
 F(S)=c_0e^{-\ii S},\qquad |c_0|=1.
\end{equation}
Normal-family extraction, including the phase at zero, supplies
such locally uniform convergence.

The normalized second field tends to zero. In fact, for each fixed
\(C_1\), use disks of radius a fixed small multiple of \(\ell_n\)
throughout \(|\Im s|\le h_n+C_1\ell_n\), with an extra margin
of a few such radii. Symmetry and
\eqref{profile:linear-log-maximum} give the global log bound
\(K=v+O_{C_1}(1)\) on their heights. Equations
\eqref{profile:periodic-deficit}--\eqref{profile:periodic-q-bound}
and \eqref{profile:selected-size-bounds} imply
\begin{equation}
 \label{profile:selected-q-smallness}
 |q(s)|^{1/2}\le
 C(\ell_n^{-1}+\ell_n^{-2})(1+v)=\exp(o(v)).
\end{equation}
Consequently \(Q_n=q/m_n^2\to0\) on every scaled compact set.

We finish by verifying the coarser bounds, including their full
range to a fixed physical depth. For \(R_b\le v'\le v\),
inverse integration and
\eqref{profile:flat-slope-selection} give
\begin{equation}
 \label{profile:coarser-inverse-bound}
 \frac{h(v)-h(v')}{\ell_n}
 =\int_{v'}^v e^{\mu(v)-\mu(z)}\,\dd z
 \le\int_{v'}^v e^{\eps_R(v-z)}\,\dd z
 \le (v-v')e^{\eps_R(v-v')}.
\end{equation}
Fix any \(\alpha_0\in(1/2,1)\). At height \(h(v')\), multiply
the row bound \(|f|/m_n\le e^{-(v-v')}\) by the
\(\alpha_0\)-power of the left side of
\eqref{profile:coarser-inverse-bound}. The result is at most
\[
 (v-v')^{\alpha_0}
 \exp\bigl(-(1-\alpha_0\eps_R)(v-v')\bigr),
\]
which is bounded uniformly. Hence the required power bound for
\(|f|/m_n\) holds globally in the horizontal variable, at
\(\ell_n\le\delta=h_n-\Im s\le h_n-h(R_b)\).
The remaining lower compact range has bounded \(f,q\).
Furthermore
\begin{equation}
 \label{profile:selected-log-range}
 \log(1+h_n/\ell_n)=o(v)
\end{equation}
by \eqref{profile:selected-size-bounds}. This absorbs that compact
range and, using \eqref{profile:selected-q-smallness}, the entire
\(q\) contribution: their normalized size is
\(e^{-v+o(v)}\), whereas the largest power of
\(\delta/\ell_n\) in question is only \(e^{o(v)}\).
We have therefore proved
\[
 M(s)/m_n\le C(\ell_n/\delta)^{\alpha_0}
 \qquad(\ell_n\le\delta\le h_n),
\]
uniformly in horizontal position. In the finite-width case retain
any fixed smaller range \(\delta\le\delta_b<H/2\). In the
entire case \eqref{profile:selected-log-range} is exactly the
additional condition \eqref{profile:entire-log-condition}.

All hypotheses of \Cref{thm:exponential-obstruction} now hold, with
the first profile in \eqref{profile:exponential-pairs}, by
\eqref{profile:selected-exponential} and
\eqref{profile:selected-q-smallness}. Its contradiction proves
both assertions.
\end{proof}

We can now assemble the growth alternatives. In the nonisotropic,
noncircular case, \Cref{thm:infinite-growth-exclusion} makes \(H\)
finite and the upper growth order finite. An upper order less than one
would give a uniform bound by a power of inverse depth with exponent
less than one, contrary to \Cref{prop:subcritical-exclusion}. A finite
upper order greater than one is excluded by
\Cref{thm:finite-superlinear-exclusion}. Thus the only remaining case
has a finite shape wall and upper growth order exactly one. This is the
critical case considered next.
